% SPDX-License-Identifier: GPL-3.0-or-later
% Copyright (C) 2026 Alessandro Greco
\chapter{Открытые проблемы и текущие проекты}\label{cap:31-futuro}
\aggiornatoal{23 сентября 2026 г.}{предложения torspec до номера 369; tor 0.4.9.13; Arti 2.6.0}

\begin{obiettivi}
К концу этой главы вы сможете:
\begin{itemize}[nosep]
  \item прослеживать процесс изменения протокола Tor и интерпретировать статус proposal;
  \item объяснять, почему постквантовая криптография требует изменений форматов сообщений, а не только алгоритмов;
  \item описывать Conflux, Walking Onions и другие проекты, направленные на производительность и масштабируемость;
  \item определять конкретные способы участия в проекте.
\end{itemize}
\end{obiettivi}

Ничто из описанного в этой книге не является окончательным. Протокол 2026 года отличается от версии 2016 года, где не было ntor-v3, управления перегрузкой и CGO; к 2036 году он снова изменится. Эта заключительная глава смотрит вперёд. Для каждой открытой проблемы мы спрашиваем: что пока работает недостаточно хорошо, что предложено и насколько далеко продвинулась реализация. Всё в следующих разделах по необходимости имеет предварительный характер: статусы предложений соответствуют репозиторию спецификаций на эталонную дату главы.

% ---------------------------------------------------------------------------
\section{Как меняется протокол}\label{sec:fut-processo}

С 2007 года изменения протокола разрабатываются через нумерованные \textbf{предложения}, в некотором смысле похожие на RFC IETF. Участник обсуждает идею в рассылке разработчиков, оформляет её в требуемом формате и, когда возникает широкое согласие, что идею стоит рассматривать, редакторы добавляют документ в репозиторий под номером предложения. Далее предложение обсуждается и уточняется, пока не становится достаточно подробным для реализации. После реализации его нормативное содержание переносится в спецификации. В отличие от RFC, предложение может изменяться, сохраняя номер, и никогда не является окончательным документом для уже реализованной функции: эту роль выполняют сами спецификации (\proposta{001}).

\begin{figure}[htbp]
\centering
% SPDX-License-Identifier: GPL-3.0-or-later
% Copyright (C) 2026 Alessandro Greco
%% Русское альтернативное описание EPUB.
\altfig{Неполная учебная схема статусов предложений Tor по предложению 001. Возможный, но необязательный путь: Draft, Open, Accepted, Finished (реализовано), Closed (включено в спецификации). Отдельно показаны Needs-Revision (требуются исправления), Needs-Research (нужны исследования), Dead (неактивно, возможно возобновление), Rejected (предложенный вариант отвергнут) и Informational (окончательный информационный документ, не являющийся нормативной спецификацией). Другие статусы, включая Meta, Reserve и Obsolete, опущены. Англоязычные названия сохранены как официальные идентификаторы.}
% Жизненный цикл Tor proposal (001-process)
\begin{tikzpicture}[font=\sffamily\scriptsize,node distance=0.55cm,
  st/.style={draw=torviola,fill=torviola!8,rounded corners=2pt,minimum width=1.6cm,minimum height=0.6cm},
  lat/.style={draw=torgrigio,fill=torgrigio!8,rounded corners=2pt,minimum width=1.6cm,minimum height=0.55cm}]
  \node[st] (d) {Draft};
  \node[st,right=of d] (o) {Open};
  \node[st,right=of o] (a) {Accepted};
  \node[st,right=of a] (f) {Finished};
  \node[st,right=of f] (c) {Closed};
  \node[above=0.25cm of a,text=torgrigio] {возможный, не обязательный путь};
  \draw[->] (d) -- (o); \draw[->] (o) -- (a); \draw[->] (a) -- (f); \draw[->] (f) -- node[above] {в спецификацию } (c);
  \node[lat,below=0.8cm of o,xshift=-1.2cm] (nr) {Needs-Revision};
  \node[lat,right=0.3cm of nr] (ns) {Needs-Research};
  \node[lat,right=0.3cm of ns] (dead) {Dead};
  \node[lat,right=0.3cm of dead] (rej) {Rejected};
  \node[below=0.22cm of ns,text=torgrigio,align=center] {другие статусы, не цепочка};
  \node[lat,below=0.85cm of nr] (inf) {Informational};
  \node[below=0.1cm of inf,text=torgrigio,align=center] {информационный документ\\не нормативная спецификация};
\end{tikzpicture}

\caption{Статусы предложений Tor согласно Proposal~001 (упрощённая и неполная схема). Путь \emph{Draft}--\emph{Open}--\emph{Accepted}--\emph{Finished}--\emph{Closed} лишь пример: \emph{Finished} означает реализацию, а \emph{Closed} также включение в спецификации. \emph{Needs-Revision} и \emph{Needs-Research} обозначают разные препятствия; \emph{Dead} может вернуться в обсуждение, тогда как \emph{Rejected} означает отказ от предложенного решения. \emph{Informational} — итоговый информационный документ, не нормативная спецификация. Другие статусы, включая \emph{Meta}, \emph{Reserve} и \emph{Obsolete}, не показаны.}
\label{fig:proposte}
\end{figure}

Статус предложения (\cref{fig:proposte}) — первое, что следует проверить читателю. В книге встречались предложения в разных состояниях. ntor-v3 (\proposta{332}) имеет статус \emph{Closed}: реализован и включён в спецификации. Управление перегрузкой (\proposta{324}) и Conflux (\proposta{329}) — \emph{Finished}: реализованы, но ещё не полностью перенесены в спецификации. CGO (\proposta{359}) остаётся \emph{Open}, хотя уже присутствует и в tor, и в Arti. План удаления TAP (\proposta{350}) имеет статус \emph{Accepted}, а Walking Onions (\proposta{300}) — \emph{Informational}.

% ---------------------------------------------------------------------------
\section{Квантовые компьютеры}\label{sec:fut-pq}

\paragraph{Проблема.} Все механизмы открытого ключа Tor, включая X25519 для согласования ключей, Ed25519 для подписей и оставшиеся применения RSA, будут нарушены достаточно мощным квантовым компьютером (\cref{cap:05-crittografia}). Никто не знает, появится ли он и когда. Однако для системы анонимности угроза может начаться раньше: атакующий способен записывать зашифрованный трафик сегодня и расшифровать его в будущем.

\paragraph{Два сценария.} \proposta{355} от 2025 года различает две ситуации. В \emph{переходной} среде достаточно защиты от \emph{будущего} квантового компьютера: записанный сегодня трафик должен выдержать последующее расшифрование, тогда как традиционные подписи остаются приемлемыми, поскольку активная подделка подписи потребовала бы квантового компьютера уже сегодня. В системе \emph{нового поколения} защита должна учитывать и квантовый компьютер, который \emph{уже существует} и способен подделывать подписи или активно атаковать согласование ключей.

\paragraph{Цена.} В 2024 году NIST стандартизировал ML-KEM, постквантовый механизм установления ключей, а также алгоритмы подписи ML-DSA и SLH-DSA~\cite{nistpqc2024}. Proposal сравнивает их размеры с современными алгоритмами:
\begin{center}\small
\begin{tabular}{@{}lrr@{}}
\toprule
Алгоритм & Открытый ключ & Шифртекст или подпись \\
\midrule
X25519 & 32 байта & 32 байта \\
ML-KEM-768 & 1\,184 байта & 1\,088 байт \\
Ed25519 & 32 байта & 64 байта \\
ML-DSA-65 & 1\,952 байта & 3\,309 байт \\
\bottomrule
\end{tabular}
\end{center}
В \cref{cap:08-celle} мы видели, что relay message может содержать не более 498 байт. Один открытый ключ ML-KEM-768 потребовал бы трёх ячеек ещё до учёта остальных данных handshake. Поэтому постквантовая криптография — не просто замена алгоритма: она меняет размещение сообщений внутри ячеек. В нынешнем формате relay нельзя непосредственно передать сообщение такого размера: три ячейки — лишь нижняя граница по объёму, если новый формат допускает фрагментацию.

\paragraph{Текущее направление.} Выводы предложения осторожны. В переходный период может быть достаточно добавить постквантовый механизм к ntor, как предлагали более ранние proposal. Для системы нового поколения предпочтительная конструкция использует два постквантовых механизма установления ключей, чтобы дополнительно обеспечить прямую секретность, если стоимость окажется приемлемой. Добавление постквантового ключа в микродескрипторы также частично съест экономию пропускной способности, достигнутую удалением TAP-ключей (\proposta{350}) и внедрением семейств, основанных на ключах (\cref{sec:relay-famiglie}). Решения зависят от ещё необходимых измерений: как часто цепи строятся с нуля и сколько пропускной способности и процессорного времени ретрансляторов уходит на согласование ключей? Это остаётся открытой исследовательской задачей и возможностью внести вклад измерениями (\cref{cap:29-ricerca}).

% ---------------------------------------------------------------------------
\section{Более крупные сообщения — и более мелкие}\label{sec:fut-340}

\proposta{340}, имеющий статус \emph{Open} на эталонную дату, напрямую решает ограничение «одно ретрансляционное сообщение на ячейку». Оно отделяет ретрансляционные сообщения от ячеек: длинные сообщения смогут занимать несколько ячеек, а несколько коротких — совместно использовать одну. Мотивация включает и сообщения длиннее 498 байт, необходимые для постквантового согласования ключей, передачи UDP и Walking Onions, и очень короткие сообщения вроде \texttt{SENDME}, \texttt{XON} и \texttt{XOFF}. Сейчас каждое из них занимает целую ячейку; новый формат позволит скрывать их в пространстве заполнения других сообщений, делая невидимыми для наблюдателя. Формат проектируется для работы с обновлённым шифрованием ретрансляционных сообщений, чья схема CGO рассмотрена в \cref{sec:cgo}; координация изменений позволит избежать множества несовместимых форматов.

% ---------------------------------------------------------------------------
\section{Производительность: несколько путей}\label{sec:fut-conflux}

Управление перегрузкой (\cref{sec:cc}) сделал отдельные цепи эффективнее. Следующий шаг — \textbf{Conflux} (\proposta{329}, \emph{Finished}). Клиент строит две цепи к одному выходному узлу с разными сторожевыми и промежуточными ретрансляторами, связывает их и распределяет трафик с учётом измеренного времени прохождения сигнала туда и обратно по каждой цепи. Идея возникла в исследованиях AlSabah, Bauer, Elahi и Goldberg; версия Tor адаптирует её, используя данные управления перегрузкой и опыт многопутевого TCP. Мы наблюдали цепи Conflux в лабораторной сети (\cref{sec:rp-stem}) и через Stem (\cref{cap:18-stem}). Предложение рекомендует, чтобы две ветви не имели общих ретрансляторов, кроме выходного узла, чтобы избегать общих узких мест. Это рекомендация, а не жёсткое правило, и на маленькой лабораторной сети она часто нарушается. tor использует Conflux с версии 0.4.8; Arti, как отмечено в \cref{sec:arti-supporto}, на эталонную дату ещё не строит такие цепи. Точки OnionPerf, работающие с Conflux и без него, позволяют сравнивать эффект в реальной сети (\cref{sec:metrics-prestazioni}).

% ---------------------------------------------------------------------------
\section{Масштабируемость: Walking Onions}\label{sec:fut-walking}

Сейчас каждый клиент Tor знает \emph{каждый} relay: он загружает консенсус и microdescriptors, поэтому трафик каталогов растёт вместе с числом ретрансляторов. Микродескрипторы и разностные документы консенсуса уменьшают стоимость, но не меняют фундаментальную форму. Начальная стоимость также сосредоточена на начальной загрузке: новый клиент должен скачать полный каталог до построения первой цепи (\proposta{300}).

У очевидных решений есть известные недостатки. Если показывать каждому клиенту лишь часть сети, возникают \emph{эпистемические} атаки: различия между представлениями сети могут идентифицировать клиентов. Если позволить каждому ретранслятору выбирать следующий, как в некоторых децентрализованных системах, вредоносные ретрансляторы получают возможность перехватывать маршруты. Walking Onions предлагает иной подход: центры каталогов публикуют подписанную структуру, связывающую каждый ретранслятор с числовым интервалом. Клиент случайно выбирает число, а предыдущий ретранслятор цепи предоставляет подписанное доказательство, указывающее ретранслятор, занимающий этот интервал, вместе с запросом расширения цепи. Клиент проверяет подпись, не получая полного списка. Исходное предложение имеет статус \emph{Informational}; более подробное продолжение (\proposta{323}) пытается закрыть оставшиеся пробелы. Более новые предложения, включая постквантовые, уже учитывают эту архитектуру.

% ---------------------------------------------------------------------------
\section{Другие открытые направления}\label{sec:fut-altri}

\begin{description}
  \item[UDP] Tor передаёт только TCP. Поэтому видеозвонки, игры и QUIC, транспорт под HTTP/3, не поддерживаются или требуют потенциально небезопасных обходных решений. \proposta{348} (\emph{Open}) рассматривает варианты от туннелей к выходным узлам до непосредственной передачи датаграмм, одновременно учитывая новые риски: отпечатки приложений, сбор метаданных однорангового обмена и исчерпание портов.
  \item[Распознавание сайтов по трафику] В конце \cref{sec:att-wf} осталась открытая проблема: универсальная защита от распознавания сайтов по трафику с приемлемыми накладными расходами пока не найдена. Автоматы заполнения, уже используемые tor для цепей onion-сервисов и разрабатываемые в Arti на основе библиотеки \emph{maybenot}, создают основу для будущих средств защиты.
  \item[Корреляция] Фундаментальная проблема \cref{cap:26-attacchi} остаётся нерешённой: сеть с низкой задержкой может сохранять коррелируемые временные и объёмные закономерности, несмотря на буферизацию и параллельный трафик. Известные защиты вроде заполнения с постоянной скоростью и дополнительных задержек дороги; исследования продолжаются в поиске более эффективных компромиссов и выбора маршрута с учётом интернет-топологии.
  \item[DoS] Доказательство выполнения работы защищает onion-сервисы (\cref{cap:15-dos}); автоматическая настройка механизма является темой \proposta[Предложения]{362}.
  \item[Arti как ретранслятор] Когда Arti сможет полноценно работать ретранслятором и центром каталогов, C-кодовая база со временем может быть выведена из эксплуатации. Это одна из долгосрочных задач Tor Project (\cref{cap:30-arti}).
\end{description}

% ---------------------------------------------------------------------------
\section{Как внести вклад}\label{sec:fut-contribuire}

У читателей, дошедших до этого места, есть несколько способов помочь, и не все требуют программирования:
\begin{itemize}
  \item \textbf{Предоставить пропускную способность}: запустить промежуточный ретранслятор или мост либо, соблюдая меры из \cref{cap:25-legale}, выходной ретранслятор (\cref{cap:22-relay}); альтернативно можно запустить прокси Snowflake даже из браузера.
  \item \textbf{Читать и обсуждать предложения}: многие остаются \emph{Open}, потому что им нужны комментарии, анализ и измерения. Рассылка разработчиков публична.
  \item \textbf{Писать код}: tor и Arti разрабатываются в GitLab Tor Project; есть задачи с метками для новичков. Arti особенно важен для новых функций.
  \item \textbf{Проводить исследования}: использовать Shadow, Chutney и данные Metrics с соблюдением рекомендаций Research Safety Board (\cref{cap:29-ricerca}).
  \item \textbf{Переводить и документировать}: добровольцы переводят Tor Browser, его руководство и сайт. Документация всегда полезна на разных языках, включая итальянский.
  \item \textbf{Использовать Tor}: большее сообщество может расширить множество анонимности, но одно лишь число пользователей не гарантирует неразличимости, если отличаются время, поведение и метаданные.
\end{itemize}

% ---------------------------------------------------------------------------
\section*{Итоги}
\begin{itemize}
  \item Протокол развивается через нумерованные предложения, статус которых движется от \emph{Draft} к \emph{Closed} по мере реализации изменений и включения их в спецификации.
  \item Постквантовая криптография важна ещё до появления мощных квантовых компьютеров; она требует сообщений больше одной ячейки и мотивирует предложение 340 о фрагментации и объединении сообщений.
  \item Conflux распределяет трафик по двум цепям; Walking Onions потенциально устраняет необходимость каждому клиенту знать всю сеть.
  \item UDP, защита от распознавания сайтов по трафику и корреляции, настройка DoS-защиты и переход Arti к роли ретранслятора остаются открытыми задачами.
\end{itemize}

% ---------------------------------------------------------------------------
\section*{Упражнения}

\begin{esercizio}[1]\label{es:31-stato}
Выберите три proposal, упомянутые в этой книге. Для каждого определите текущий статус и его значение: развёрнута ли функция и описана ли она в спецификациях?
\end{esercizio}

\begin{esercizio}[1]\label{es:31-celle}
Сколько ячеек требуется для передачи одного открытого ключа ML-KEM-768 и одного шифртекста ML-KEM-768, если одно ретрансляционное сообщение содержит максимум 498 байт, а отдельные сообщения не могут совместно использовать одну ячейку? Что изменится при использовании \proposta[Предложения]{340}? Различите теоретическую нижнюю границу числа ячеек (без служебных данных) и фактическое число, зависящее от упаковки и фрагментации по предложению.
\end{esercizio}

\begin{esercizio}[2]\label{es:31-epistemico}
Проиллюстрируйте epistemic-атаку: если разные клиенты знают разные подмножества сети, как атакующий может использовать различия?
\end{esercizio}

\begin{esercizio}[2]\label{es:31-conflux}
Используя Stem в лабораторной сети, найдите пары Conflux circuits с общим \file{CONFLUX\_ID}. Проверьте, заканчиваются ли обе ветви на одном exit и избегают ли остальных общих relay, как рекомендует \proposta{329}. Что вы наблюдаете и как это объяснить?
\end{esercizio}

\begin{esercizio}[3]\label{es:31-proposta}
Выберите открытую проблему из этой главы и напишите короткий proposal в формате Tor Project: мотивация, дизайн, анализ безопасности, совместимость и план реализации.
\end{esercizio}

% ---------------------------------------------------------------------------
\section*{Дополнительная литература}
\begin{itemize}
  \item Индекс proposal в репозитории \texttt{torspec}, организованный по статусам, и документ процесса proposal (\proposta{001}).
  \item N.~Mathewson, \titolo{Options for postquantum circuit extension handshakes} (\proposta{355}).
  \item \titolo{Walking Onions: Scaling and Saving Bandwidth} (\proposta{300}) и продолжение (\proposta{323}).
  \item Блог Tor Project, где объявляются и объясняются изменения протокола.
\end{itemize}
